\section{Introduction}
\label{sec:introduction}

Large language models (LLMs) have recently been used as reasoning engines in financial analysis. Their value for knowledge-based decision support depends on how quantitative, textual and visual evidence is represented and assessed, rather than on distributing tasks among agents alone. Frameworks such as QuantAgent \cite{xiong2025quantagent} provide a starting point for chart-based LLM decision support. Adapting such a framework to Vietnamese equities requires explicit treatment of Vietnamese-language news, single-asset factor inputs and the distinction between a directional forecast and a feasible cash-equity trade.

Vietnam motivates this implementation in three ways. First, the bar-count forecast horizon must be distinguished from entry, exit and settlement assumptions. Second, Vietnamese news requires a language-appropriate scorer, whose financial-domain validity still needs independent assessment. Third, OHLCV-only factor construction cannot retain the peer-relative meaning of every cross-sectional expression. These are practical adaptation questions, not evidence that the resulting system has solved emerging-market prediction.

We examine a concrete multi-agent implementation for Vietnamese equities. It combines deterministic technical-indicator reports, a native selector over 85 OHLCV factor candidates, Vietnamese sentiment scores, two chart-reading roles and a final binary Decision Agent. Indicator, Pattern, Trend and Decision roles and associated utilities reuse or modify QuantAgent components \cite{xiong2025quantagent}. The integration is an implementation and domain-adaptation exercise: condensed reports enter a prompt with heuristic evidence-priority instructions, rather than a learned knowledge-fusion representation.

Our question is how this implementation compares with class-aware and numerical controls on the same recorded historical market units, and what its archived executions can support about variability, components and economics. The primary comparison is Full versus Base. Base removes both selected-factor and sentiment sections, so this is a bundle contrast, not an isolated alpha-information treatment. A separate factorial panel includes factor-only, sentiment-only and matched-structure past-donor variants. Different prompt families and small repeated panels remain explicit rather than being pooled into a larger apparent experiment.

The completed release contains 102 main market units spanning six historical blocks and twelve selected dates, with 25 UP and 77 DOWN net labels. Full has pooled balanced accuracy (BA) of \FullBA\%, Base \BaseBA\%, and a price-only logistic comparator \LogisticPriceBA\%. CASH has higher ordinary accuracy than Full but BA of 50\%. Mean reference account returns are negative for Full and Base. These observations motivate an assessment of the implementation and its limits, rather than a claim of profitable trading or general predictive superiority.

The evaluation also exposes unresolved methodological defects. Native factor operations are isolated from prices after the outer decision cutoff in synthetic probes, but are not causal row by row within their trailing history. Native ranking and orientation diagnostics are in-sample. News availability and the sentiment classifier's financial-domain validity are not independently authenticated. We preserve these facts when reporting the frozen outputs. Separate train-fit numeric controls use causal adapters and cannot retrospectively repair the native LLM study.

The main contributions are as follows:
\begin{itemize}
\item A Vietnamese-equity adaptation that integrates an OHLCV factor selector and local-language sentiment with Indicator, Pattern, Trend and Decision roles. It extends reused chart-agent components through explicit quantitative and textual inputs rather than claiming a new learned fusion architecture.
\item A specified single-asset signal-construction method: 80 retained Formulaic Alpha adaptations and five custom handlers, top-five diagnostic ranking, and a separately represented time-decayed sentiment report. Its operator semantics and construction limits make the actual inputs distinguishable from idealised source formulas.
\item A shared-report evaluation design combining common-support baselines, repeated upstream and Decision executions, and five variants for factors, sentiment and a past donor. The comparison identifies which inputs differ without treating a bundled contrast as a causal alpha effect.
\item An empirical analysis linking class-aware prediction to transaction-cost-aware long-only accounts. On the main support, the Full--Base balanced-accuracy difference accompanies fewer false LONG forecasts and unchanged UP recall, while mean net returns remain negative. This distinguishes a classification change from demonstrated economic value.
\end{itemize}

Historical snapshots, actual requests/responses and metrics are retained with hashes for offline replay. The evidence is historical and exploratory: the reduced scope was amended after some inference, all outcome dates preceded the freeze, and the study does not establish causal component value or prospective deployment performance.

The remainder of this paper is organised as follows. Section~\ref{sec:relatedworks} reviews related work; Section~\ref{sec:formulation} formalises the task and signals; Section~\ref{sec:systemarchitecture} describes each agent; Section~\ref{sec:experimentalsetup} specifies the evaluation; Section~\ref{sec:backtestresults} reports results; Section~\ref{sec:limitations} discusses limitations and future work; and Section~\ref{sec:conclusion} concludes.
